% part_ii_posthoc.tex — Part II: post-hoc validity of (1) at θ̂
% Includes Sol-derived set-valued ZJ transfer (2026-07-27_setvalued_thm2.txt)
\documentclass[11pt]{article}
\usepackage{amsmath,amssymb,amsthm}
\usepackage[margin=1in]{geometry}
\newtheorem{definition}{Definition}
\newtheorem{assumption}{Assumption}
\newtheorem{theorem}{Theorem}
\newtheorem{corollary}{Corollary}
\newtheorem{remark}{Remark}
\newtheorem{problem}{Problem}

\title{Part II --- Post-hoc validity of $\mathcal{C}$ at $\hat\theta(\mathcal{D})$}
\author{research/formal (working)}
\date{27 July 2026}

\begin{document}
\maketitle

\noindent\textit{Depends on \texttt{common\_setup.tex} and Part~I.
Goal (iii): data-randomize selection; \textbf{never} recalibrate $\mathcal{C}$.
Set-valued transfer: \texttt{proofs/sol\_gaps/2026-07-27\_setvalued\_thm2.txt}.}

\section{Targets}

\begin{definition}[Failure indicator]
For a family $C^{(q)}(X;D,\theta)$ calibrated at miscoverage level $q$,
\[
\ell_q(W,\theta):=\mathbf{1}\{Y\notin C^{(q)}(X;D,\theta)\},
\quad W=(D,X,Y).
\]
Writeup (1) is the special case $q=\alpha$ for \emph{prespecified} $\theta$.
\end{definition}

\section{Goal (i) --- set-valued stability transfer}

\begin{assumption}[Joint fixed-$\theta$ validity on the selection support]
\label{ass:FV}
Let $\Theta_{\mathrm{val}}\subseteq\Theta$ be the (problem-specific) set of hyperparameters
for which writeup~(1) holds at miscoverage level $q$. Assume
\[
\sup_{\theta\in\Theta_{\mathrm{val}}}
\mathbb{P}\bigl\{Y\notin C^{(q)}(X;D,\theta)\bigr\}\le q.
\tag{FV}
\]
Probability is over the joint experiment $(D,X,Y)$ (and set randomization);
$\theta\in\Theta_{\mathrm{val}}$ is held fixed.
\textit{Notation:} $C^{(q)}(X;D,\theta)$ is writeup~$\mathcal{C}(X;\mathcal{D},\theta)$
when the classical miscoverage level in~(1) equals $q$ (superscript $=$ level;
third argument $=$ hyperparameter $\theta$). Do not read the superscript as replacing $\theta$.
\textit{Score-threshold families:} $\Theta_{\mathrm{val}}=\{\theta:F(\theta)\ge 1-q\}$;
do not claim~(FV) on all of $[0,1]$.
\end{assumption}

\begin{theorem}[Set-valued stability transfer]
\label{thm:set-valued}
Let $\hat\theta=A(D)$ be $(\eta,\tau,\nu)$-stable with oracle $A_0\perp W$ and
good event $G$ of probability $\ge 1-\nu$ such that on $G$,
$K_w(B)\le e^\eta Q_0(B)+\tau$ for all measurable $B\subseteq\Theta$
(Zrnic--Jordan Def.~2 / (ST)).
Assume oracle validity
$\mathbb{P}\{Y\notin C^{(q)}(X;D,A_0)\}\le q$
(implied by~(FV) if $A_0\perp W$ and $A_0$ is supported on $\Theta_{\mathrm{val}}$;
\textbf{(OV) alone suffices} ---~(FV) is only a convenient sufficient condition).
Assume $C$'s downstream randomization is a common Markov kernel given $(W,\theta)$
(no unaccounted coupling of $A$ into the construction of $C$ beyond plugging $\theta$).
Then
\[
\mathbb{P}\bigl\{Y\notin C^{(q)}(X;D,\hat\theta)\bigr\}
\le e^\eta q+\tau+\nu.
\]
In particular, with $q=\delta e^{-\eta}$,
\[
\mathbb{P}\bigl\{Y\notin C^{(\delta e^{-\eta})}(X;D,\hat\theta)\bigr\}
\le \delta+\tau+\nu.
\]
\end{theorem}

\begin{remark}[Proof idea]
Failure section $B_q(w)=\{\theta:y\notin C^{(q)}(x;d,\theta)\}$; apply (ST) to the
data-dependent section; add $\nu$-mass of $G^c$; use $A_0\perp W$.
\end{remark}

\begin{remark}[Goals (i)--(ii) vs (iii)]
Using $C^{(\delta e^{-\eta})}$ instead of $C^{(\alpha)}$ is classical-level inflation
(recalibration of the \emph{level}). That addresses (i)--(ii) when $\eta$ is fixed.
Goal~(iii) forbids rebuilding $\mathcal{C}$; it drives $\eta,\tau,\nu\to 0$ by
randomizing \emph{selection} only.
\end{remark}

\section{Goal (ii) --- inflation}

\begin{definition}[Inflation]
$\mathrm{Infl}:=\mathbb{P}\{Y\notin C^{(\alpha)}(X;D,\hat\theta)\}-\alpha$.
Theorem~\ref{thm:set-valued} implies
$\mathrm{Infl}\le (e^\eta-1)\alpha+\tau+\nu$ when $q=\alpha$ is used without level change.
\end{definition}

\section{Goal (iii) --- vanishing stability correction}

\begin{corollary}[Vanishing inflation at fixed classical level --- W5/(iii)]
\label{cor:vanish}
Let $\{A_r\}$ be $(\eta_r,\tau_r,\nu_r)$-stable with
$\eta_r\to 0$ and $\tau_r+\nu_r\to 0$, and suppose randomizing $A_r$ does
\textbf{not} alter the fixed-$\theta$ map $(D,X)\mapsto C^{(\alpha)}(X;D,\theta)$
(same classical level $\alpha$ as in writeup~(1); \emph{no} level change).
If~(OV) holds for the oracle of $A_r$ at level $q=\alpha$ and the fixed-$\theta$ map
is unchanged, then Theorem~\ref{thm:set-valued} with
$q=\alpha$ yields
\[
\mathbb{P}\bigl\{Y\notin C^{(\alpha)}(X;D,A_r(D))\bigr\}
\le e^{\eta_r}\alpha+\tau_r+\nu_r,
\]
hence
\[
\mathrm{Infl}_r
\le (e^{\eta_r}-1)\alpha+\tau_r+\nu_r
\to 0.
\]
In the limit $\eta=\tau=\nu=0$, classical (1) holds at the selected parameter
with the \emph{same} $C^{(\alpha)}$.
(The optional path $C^{(\delta e^{-\eta})}$ belongs only under goals~(i)--(ii);
it is \textbf{not} the (iii) mechanism.)
\end{corollary}

\begin{definition}[Randomized selection map]
$(\mathcal{D},\xi)\mapsto\tilde\theta(\mathcal{D};\xi)$ with $\xi$ independent of the
test point; fixed-$\theta$ map of $\mathcal{C}$ unchanged.
\end{definition}

\begin{corollary}[Forbidden repairs]
Sample-splitting re-estimation of $\mathcal{C}$ and threshold/score recalibration of
$\mathcal{C}$ are out of scope as answers to (iii) (contrast literature only).
\end{corollary}

\begin{problem}[Remaining]
Instantiate $A_r$ via \texttt{part\_i\_randomized\_design.tex}
(RNM / Lap $b=2\varepsilon/\eta_r$, or soft-argmin $\beta=\eta_r/(2\varepsilon)$)
with $\eta_r\to 0$; verify (FV) and Ass.\ concentration $\varepsilon(n),\nu(n)$ for the
project's $(\hat f,\hat g)$ and $\mathcal{C}$.
\end{problem}

\end{document}
